\section{Riscos, conformidade e limites}
Os riscos prioritários são: (a) risco de confiança, caso a triagem seja interpretada como diagnóstico; (b) risco financeiro do fundo de créditos; (c) risco de liquidez e repasse às clínicas; (d) risco de privacidade e segurança; (e) risco operacional de atendimento manual; (f) risco de baixa disposição a pagar; e (g) risco de métricas infladas por downloads, cadastros incompletos ou amostras de conveniência. As salvaguardas incluem revisão veterinária, linguagem não diagnóstica, consentimento informado, minimização de dados, trilha de auditoria, contratos claros, limites de créditos e divulgação transparente de resultados.

As evidências disponíveis são exploratórias e documentais. Não foi identificada execução do piloto de oito semanas, contrato de clínica, receita efetivamente recebida ou medição real de CAC/retenção. O plano deve ser atualizado após a coleta. A amostra de conveniência não permite generalizar resultados para todos os tutores de Belo Horizonte, muito menos para o Brasil.